\documentclass[10pt,a4paper]{amsart}
\usepackage[margin=19mm]{geometry}
\usepackage{amsmath,amssymb,mathtools}
\usepackage[scr=rsfs,cal=euler]{mathalfa}
\usepackage{xfrac}
\usepackage[unicode,hidelinks,bookmarks=false]{hyperref}
\newcommand{\ie}{i.e.\ }
\newcommand{\cf}{cf.\ }
\newcommand{\resp}{respectively\ }
\newcommand{\Subsection}{\subsection}
\providecommand{\nobreakdash}{\nobreak\hbox{-}}
\setlength{\emergencystretch}{2em}
\multlinegap0pt
\usepackage{amssymb}
\usepackage{float}
\usepackage{tikz}
\newtheorem{theorem}{Theorem}
\title{Helly's Theorem--A Very Early Introduction}
\author{Eric L. Grinberg}
\date{Source: arXiv:2604.01237v1 (CC BY 4.0)}
\begin{document}
\maketitle

\noindent\emph{Source and licence.} Helly's Theorem--A Very Early Introduction. Version arXiv:2604.01237v1 (CC BY 4.0), distributed by arXiv under the Creative Commons Attribution 4.0 International licence, http://creativecommons.org/licenses/by/4.0/, as recorded in the arXiv licence metadata for this exact version and shown on the abstract page https://arxiv.org/abs/2604.01237v1. Full licence notice: \texttt{licenses/arxiv-2604.01237-CC-BY-4.0.txt}.\par\smallskip

\noindent\textit{Editorial scope.} Counted unit: the article's theorem theorem (source0.tex lines 150--152). The article's complete original proof of it is delivered with it (source0.tex lines 154--159, one \texttt{proof} environment of 6 lines). No mathematics is written, repaired or re-derived: the statement and the proof are the article's own lines, in the article's own notation, and no retained line is substituted or re-flowed.  No further source-local context is needed: every cross-reference of every retained line resolves inside the two delivered spans.  Nothing was omitted from any retained span, so the package carries no omission record.  No retained line contains a citation, so the wrapper carries no bibliography and no bibliography entry.  No retained line contains a figure environment, an image-inclusion command or drawing code, so no image is delivered and the package declares no asset dependency.  Editorial formatting only, in addition: an AMS \texttt{amsart} preamble that restates the article's own declarations and macros used by the retained lines, namely the environments \texttt{theorem} on separate counters, together with the standard packages those lines need.  The retained environments print in this document as Theorem 1 in order of appearance, whereas the article prints the same items with the numbering of its own full document.  The complete native source of the article as distributed in the arXiv e-print for this exact version is bundled unchanged as \texttt{sources/197645/source0.tex}.
\begin{theorem}[Helly for Planes in $\mathbb R^3$] \, \\
Let $T$ be a nondegenerate  system of $N$ linear equations in $3$ unknowns, with $N \ge 4$. Assume that any subsystem of $T$ comprised of $4$ equations is consistent. Then the full system is consistent.
\end{theorem}

\begin{proof}By nondegeneracy, each equation in the system $T$ has a solution set that is a plane. If all $N$ planes are the same, the solution set of $T$ is that plane, and the system is consistent. 

Otherwise, there are two equations, $E_1,E_2$, whose solution sets are distinct planes $P_1,P_2$. These planes must meet since we can append two additional equations to $E_1,E_2$ to form a consistent $4$-equation subsystem, and solutions to that subsystem belong to both planes. Thus $P_1 \cap P_2$ is a line $\ell$, since two distinct intersecting planes in $3$-space meet in a line.

If all the rest of the solution planes meet $P_1 \cap P_2$  in $\ell$ then $\ell$ is part of the solution set of $T$ and hence $T$ is consistent.  If some solution plane $P_3$ meets $P_1 \cap P_2$ in less than the full line, that solution plane meets $\ell$ at just one point $Q$. Each additional solution plane beyond $P_1,P_2,P_3$  must also meet $\ell$ at $Q$, else the resulting $4$-equation subsystem would be inconsistent. Hence the point $Q$ belongs to all the solution planes and so $T$ is consistent.
\end{proof}

\end{document}
